\documentclass[11pt]{article}

\usepackage[margin=1in]{geometry}
\usepackage[T1]{fontenc}
\usepackage{lmodern}
\usepackage{microtype}
\usepackage{amsmath,amssymb,mathtools,amsthm}
\usepackage{booktabs}
\usepackage{array}
\usepackage{enumitem}
\usepackage{xcolor}
\usepackage{hyperref}
\usepackage[nameinlink,noabbrev]{cleveref}
\usepackage{listings}

\hypersetup{
  colorlinks=true,
  linkcolor=blue!55!black,
  citecolor=blue!55!black,
  urlcolor=blue!55!black,
  pdftitle={A resolvent-algebra version of Algorithm 5.1},
  pdfauthor={OpenAI}
}

\lstset{
  basicstyle=\ttfamily\small,
  columns=fullflexible,
  frame=single,
  breaklines=true,
  xleftmargin=0.5em,
  xrightmargin=0.5em,
  aboveskip=0.7em,
  belowskip=0.7em
}

\setlist[enumerate]{itemsep=0.35em,topsep=0.45em}
\setlist[itemize]{itemsep=0.25em,topsep=0.4em}
\setlength{\parindent}{0pt}
\setlength{\parskip}{0.55em}

\newcommand{\QQ}{\mathbf{Q}}
\newcommand{\ZZ}{\mathbf{Z}}
\newcommand{\PP}{\mathbf{P}}
\newcommand{\Aff}{\mathbf{A}}
\newcommand{\GG}{\mathbf{G}}
\newcommand{\ol}[1]{\overline{#1}}
\newcommand{\Gal}{\operatorname{Gal}}
\newcommand{\Pic}{\operatorname{Pic}}
\newcommand{\Spec}{\operatorname{Spec}}
\newcommand{\Tr}{\operatorname{Tr}}
\newcommand{\Nm}{\operatorname{Nm}}
\newcommand{\Hom}{\operatorname{Hom}}
\newcommand{\End}{\operatorname{End}}
\newcommand{\Mat}{\operatorname{Mat}}
\newcommand{\Span}{\operatorname{Span}}
\newcommand{\im}{\operatorname{im}}
\newcommand{\rank}{\operatorname{rank}}
\newcommand{\Fix}{\operatorname{Fix}}
\newcommand{\Stab}{\operatorname{Stab}}
\newcommand{\disc}{\operatorname{disc}}
\newcommand{\PGL}{\operatorname{PGL}}
\newcommand{\cM}{\mathcal{M}}
\newcommand{\tcM}{\widetilde{\mathcal{M}}}
\newcommand{\tM}{\widetilde{M}}
\newcommand{\cA}{\mathcal{A}}
\newcommand{\cO}{\mathcal{O}}
\newcommand{\cP}{\mathcal{P}}
\newcommand{\Gamhat}{\widehat{\Gamma}}
\newcommand{\gamhat}{\widehat{\gamma}}
\newcommand{\Deltav}{\Delta_{\mathrm{VdM}}}
\newcommand{\deltaE}{\delta_{E_6}}
\newcommand{\DeltaE}{\Delta_{E_6}}
\newcommand{\W}{W(E_6)}
\newcommand{\V}{V_6}
\newcommand{\QOmega}{\QQ^{\Omega}}
\newcommand{\red}[1]{\textcolor{red!75!black}{#1}}

\newtheorem{definition}{Definition}[section]
\newtheorem{proposition}[definition]{Proposition}
\newtheorem{lemma}[definition]{Lemma}
\newtheorem{theorem}[definition]{Theorem}
\newtheorem{remark}[definition]{Remark}

\newenvironment{summarybox}{%
  \begin{center}\begin{minipage}{0.94\textwidth}\small
  \hrule\vspace{0.55em}%
}{%
  \vspace{0.55em}\hrule\end{minipage}\end{center}%
}

\newenvironment{implementationbox}{%
  \begin{center}\begin{minipage}{0.94\textwidth}\small
  \hrule\vspace{0.55em}\textbf{Implementation point.}\ }{%
  \vspace{0.55em}\hrule\end{minipage}\end{center}%
}

\title{A Resolvent-Algebra Version of Algorithm 5.1\\
\large Avoiding a full explicit $\QQ$-basis for the splitting field}
\author{Self-contained notes on the improved cubic-surface algorithm}
\date{\today}

\begin{document}
\maketitle

\begin{abstract}
Algorithm 5.1 of Elsenhans--Jahnel constructs a smooth cubic surface over
$\QQ$ with prescribed Galois action on its $27$ lines.  In the direct
implementation, one represents the splitting field $L/\QQ$ as a
$\QQ$-vector space, writes the Galois automorphisms as large rational
matrices, and solves semilinear descent equations in dimensions $6$ and
$10$.  This note describes a more economical version.  The ten-dimensional
projective descent is unnecessary if one evaluates Clebsch invariants directly
on the twisted six-dimensional source.  The remaining six-dimensional descent
can be modeled using lower-degree resolvent algebras attached to suitable
$G$-sets, especially the $27$-line $G$-set, instead of using the full splitting
field.  The resulting algorithm takes as arithmetic input a separable
polynomial $f(x)$ with labelled Galois group and a homomorphism into
$W(E_6)$; it avoids constructing a $\QQ$-basis of the full splitting field
unless one explicitly chooses the regular $G$-set.
\end{abstract}

\tableofcontents

\section{Purpose and main conclusion}

The goal is to construct a smooth cubic surface $X/\QQ$ whose induced Galois
action on
\[
  \Pic(X_{\ol\QQ})
\]
--- equivalently, on the configuration of the $27$ geometric lines --- realizes a
prescribed representation
\[
  \rho:G_{\QQ}\longrightarrow \W.
\]
Computationally, one usually replaces $G_{\QQ}$ by a finite quotient.  Thus the
input is a finite Galois extension $L/\QQ$, a finite group
$G=\Gal(L/\QQ)$, and a homomorphism
\[
  \rho:G\longrightarrow \W.
\]
The original descent-heavy implementation represents $L$ explicitly as a
$\QQ$-vector space.  This is often the dominant cost.

The improvement in this note has two parts.

\begin{enumerate}[label=\textbf{Improvement \arabic*.}]
\item Do not construct the descended ten-dimensional projective model
$\tM_\rho\subset \PP^9_\QQ$ unless one specifically wants it.  For the purpose of
producing a cubic equation, it is enough to evaluate the $W(E_6)$-invariant
Clebsch quantities directly on the twisted six-dimensional source.
\item Do not construct the full splitting field $L$ as a $\QQ$-vector space when
building the twisted six-dimensional source.  Instead use a lower-degree
resolvent algebra associated with a $G$-set $\Omega$ whose permutation module
contains the six-dimensional $E_6$ reflection representation.
\end{enumerate}

The resulting mathematical input is not merely the abstract group $G$.  One needs
both:
\[
  \boxed{\text{a polynomial }f(x)\in\QQ[x]\text{ whose splitting field is }L,}
\]
and
\[
  \boxed{\text{a labelled Galois group }G\leq S_{\deg f}\text{ together with }
  \rho:G\to\W.}
\]
The polynomial encodes the arithmetic $G$-torsor.  The labelled permutation group
records how this torsor acts on the roots.  The homomorphism $\rho$ transfers this
arithmetic torsor to the $E_6$ marking data.

\begin{summarybox}
\textbf{Main conclusion.}
The full splitting field $L$ is still present mathematically, but it need not be
represented as a $\QQ$-vector space.  It can be replaced by smaller associated
\'{e}tale algebras
\[
  \cA_\Omega=(L\otimes_\QQ \QQ^\Omega)^G
\]
attached to carefully chosen $G$-sets $\Omega$.  The universal safe choice is the
$27$-line $G$-set obtained by restricting the standard $W(E_6)$-action on the
$27$ lines along $\rho$.
\end{summarybox}

\section{The mathematical problem}

\subsection{Input}

The improved algorithm has the following input.

\begin{enumerate}[label=\textbf{Input \arabic*.}]
\item A separable polynomial
\[
  f(x)\in \QQ[x]
\]
with splitting field $L$.

\item A labelled Galois group computation: a subgroup
\[
  G\leq S_d,\qquad d=\deg f,
\]
identified with $\Gal(L/\QQ)$ through its action on an ordered list of the roots
\[
  \alpha_1,\ldots,\alpha_d.
\]
Equivalently, for each $g\in G$, the action is
\[
  g(\alpha_i)=\alpha_{g(i)}.
\]

\item A homomorphism
\[
  \rho:G\longrightarrow \W.
\]
The embedding of the image into $\W$ is part of the input; the abstract isomorphism
type of the image is not enough.

\item Universal precomputed $E_6$ data:
\begin{itemize}
\item the six-dimensional rational representation
$\V$ of $\W$ coming from Pinkham's cusp-parameter construction;
\item the $40$ normalized Coble gamma labels and their convention sign mask;
\item the formulas for the $W(E_6)$-invariant Clebsch quantities
$A,B,C,D,E$;
\item the $27$-line $W(E_6)$-set and the $W(E_6)$-equivariant map from its
permutation module onto $\V$.
\end{itemize}
\end{enumerate}

\subsection{Output}

The desired output is a homogeneous cubic form
\[
  F(X_0,X_1,X_2,X_3)\in\QQ[X_0,X_1,X_2,X_3]
\]
such that the surface
\[
  X=\{F=0\}\subset\PP^3_\QQ
\]
is smooth and the induced Galois representation on the $27$ geometric lines of
$X$ agrees with $\rho$, up to the marking convention used to identify the line
configuration with the standard $E_6$ configuration.

\subsection{Why an abstract group is insufficient}

Two different polynomials can have labelled Galois groups abstractly isomorphic to
the same subgroup of $\W$ but have non-isomorphic splitting fields.  They define
different classes in
\[
  H^1(\QQ,G),
\]
and hence can define different twists of the marked moduli space.  The abstract
group determines the possible symmetry type; the polynomial determines the actual
arithmetic torsor.

Thus the data
\[
  G\quad\text{and}\quad \rho:G\to\W
\]
are not enough.  One also needs the polynomial $f$, or an equivalent presentation
of the same $G$-torsor over $\QQ$.

\section{Universal geometry used by the algorithm}

\subsection{The marked cubic-surface moduli space}

Let $\tcM$ denote the moduli space of marked smooth cubic surfaces.  A marking is
an ordered six-tuple of pairwise disjoint lines.  After blowing down these six
lines, a marked cubic surface over an algebraically closed field is the same as an
ordered configuration of six points
\[
  p_1,\ldots,p_6\in\PP^2
\]
in general position, modulo the diagonal action of $\PGL_3$.

The Weyl group $\W$ acts on the marking.  The quotient by this action is the
unmarked moduli space.  Algorithm 5.1 prescribes a Galois action on the marking,
then forgets the marking after computing invariant moduli data.

\subsection{The $27$ lines and the six-dimensional representation}

Use the standard blow-up notation
\[
  \Pic(X_{\ol\QQ})=\ZZ H\oplus \ZZ E_1\oplus\cdots\oplus\ZZ E_6,
\]
with intersection form
\[
  H^2=1,
  \qquad
  E_i^2=-1,
  \qquad
  H\cdot E_i=0,
  \qquad
  E_i\cdot E_j=0\quad(i\neq j).
\]
The canonical class is
\[
  K=-3H+E_1+\cdots+E_6.
\]
The $27$ line classes are
\[
  E_i,\qquad
  H-E_i-E_j\quad(1\leq i<j\leq6),
\]
and
\[
  2H-\sum_{j\neq i}E_j\quad(1\leq i\leq6).
\]
Let
\[
  \V=K^\perp\otimes_\ZZ\QQ.
\]
This is the rational six-dimensional reflection representation of $\W$.

Let $\Omega_{27}$ be the set of the $27$ line classes.  There is a
$W(E_6)$-equivariant surjection
\begin{equation}
  q_{27}:\QQ^{\Omega_{27}}\longrightarrow \V
  \label{eq:q27}
\end{equation}
defined on basis vectors by
\begin{equation}
  e_\ell\longmapsto \ell+\frac{1}{3}K.
  \label{eq:line-projection}
\end{equation}
Indeed,
\[
  \left(\ell+\frac{1}{3}K\right)\cdot K
  =\ell\cdot K+\frac{1}{3}K^2
  =-1+1=0.
\]
The images of the $27$ line classes span $K^\perp\otimes\QQ$.

\begin{implementationbox}
The map $q_{27}$ is the universal reason that the $27$-line resolvent is enough:
after restricting the $27$-line action along $\rho:G\to W(E_6)$, the associated
permutation module still surjects onto the desired six-dimensional source
representation.
\end{implementationbox}

\subsection{Pinkham's cuspidal-cubic source}

Let
\[
  C:yz^2=x^3\subset\PP^2
\]
be the cuspidal cubic.  Its smooth locus is parametrized by
\[
  p(t)=(t:t^3:1),\qquad t\in\Aff^1.
\]
For
\[
  t=(t_1,\ldots,t_6)\in\Aff^6,
\]
put
\[
  p_i(t)=p(t_i).
\]
When these six points are in general position, blowing them up gives a marked
smooth cubic surface.  Thus one obtains a rational map
\[
  \Aff^6\dashrightarrow \tcM.
\]
The source $\Aff^6$ carries Pinkham's six-dimensional $W(E_6)$-action, and the map
to $\tcM$ is $W(E_6)$-equivariant.

The general-position locus is the complement of the $E_6$ root hyperplanes.  More
explicitly, define
\begin{equation}
 \boxed{
  \deltaE(t)=
  \left(\prod_{1\leq i<j\leq6}(t_j-t_i)\right)
  \left(\prod_{1\leq i<j<k\leq6}(t_i+t_j+t_k)\right)
  (t_1+\cdots+t_6).
 }
 \label{eq:deltaE}
\end{equation}
Then the regular locus is
\[
  V^\circ=\{t\in\Aff^6:\deltaE(t)\neq0\}.
\]
The factors have the following meanings:
\begin{itemize}
\item $t_i=t_j$: two of the six points coincide;
\item $t_i+t_j+t_k=0$: three of the six points are collinear;
\item $t_1+\cdots+t_6=0$: all six points lie on a conic.
\end{itemize}

\begin{summarybox}
\textbf{Important search-domain rule.}
One should not exclude $t_i=0$.  The point $p(0)=(0:0:1)$ is smooth on the
cuspidal cubic.  The only geometric exclusions in the source are the factors of
$\deltaE$ in \eqref{eq:deltaE}, plus later Clebsch, pentahedral, and smoothness
failure loci.
\end{summarybox}

For invariant boundary testing it is often preferable to use
\begin{equation}
  \DeltaE(t)=\deltaE(t)^2.
  \label{eq:DeltaE}
\end{equation}
The polynomial $\DeltaE$ is $W(E_6)$-invariant, whereas $\deltaE$ is a relative
invariant.

\section{Coble coordinates and Clebsch invariants}

\subsection{Normalized Coble coordinates on the cusp}

For $I\subseteq\{1,\ldots,6\}$, write
\[
  s_I(t)=\sum_{i\in I}t_i,
  \qquad
  s_{[6]}(t)=t_1+\cdots+t_6,
\]
and
\[
  \Delta_I(t)=\prod_{\substack{i<j\\ i,j\in I}}(t_j-t_i),
  \qquad
  \Deltav(t)=\Delta_{\{1,\ldots,6\}}(t).
\]
On the cuspidal family, every raw Coble gamma coordinate has a common
Vandermonde factor $\Deltav$.  Dividing by this factor gives normalized degree-nine
coordinates.

For a decomposition
\[
  I\sqcup I^c=\{1,\ldots,6\},\qquad |I|=|I^c|=3,
\]
the type-I normalized coordinate is
\begin{equation}
 \boxed{
  \gamhat_{I\mid I^c}(t)
  =\Delta_I(t)\Delta_{I^c}(t)s_I(t)s_{I^c}(t)s_{[6]}(t).
 }
 \label{eq:gammahatI}
\end{equation}
There are $10$ such coordinates.

For a partition of $\{1,\ldots,6\}$ into three unordered pairs $A,B,C$, and a
cyclic ordering $(A,B,C)$ of the three pairs, define
\[
  d_A=t_{a_2}-t_{a_1}\quad\text{for }A=\{a_1<a_2\},
\]
and similarly for $B$ and $C$.  The type-II normalized coordinate is
\begin{equation}
\boxed{
\begin{aligned}
 \gamhat_{A\to B\to C}(t)
 &=d_A d_B d_C
 \prod_{a\in A}s_{\{a\}\cup B}
 \prod_{b\in B}s_{\{b\}\cup C}
 \prod_{c\in C}s_{\{c\}\cup A}.
\end{aligned}}
\label{eq:gammahatII}
\end{equation}
There are $30$ such coordinates, because reversing the cyclic orientation gives
the other coordinate associated with the same three-pair partition.

Choose once and for all an ordering of the $40$ labels and a sign mask
\[
  \eta=(\eta_1,\ldots,\eta_{40})\in\{\pm1\}^{40}
\]
matching the implementation convention for Coble's determinants.  The stored
normalized gamma vector is
\[
  \Gamhat(t)=\bigl(\eta_1\gamhat_1(t),\ldots,
  \eta_{40}\gamhat_{40}(t)\bigr)^t.
\]
With the correct sign mask, this vector is exactly equivariant for a signed
permutation representation of $W(E_6)$:
\begin{equation}
  \Gamhat(M_wt)=\widehat P_w\Gamhat(t),
  \qquad w\in W(E_6),
  \label{eq:gamma-equivariance}
\end{equation}
where $M_w$ denotes Pinkham's six-dimensional matrix and $\widehat P_w$ is a signed
permutation matrix on the $40$ normalized gamma labels.

\begin{implementationbox}
The sign mask is not cosmetic.  A mismatch here usually appears later as a failure
of rationality for supposedly descended quantities.  The even power sums below are
independent of changing the signs of individual gamma coordinates, but equivariance
checks and comparisons with stored Coble data are not.
\end{implementationbox}

\subsection{Clebsch invariants}

Let
\[
  P_k(t)=\sum_{a=1}^{40}\gamhat_a(t)^k.
\]
The sum is over one representative from each signed pair of Coble coordinates.  In
the formulas below only even powers appear, so the chosen signs of the forty
coordinates do not affect the value.  One should not sum over all eighty signed
coordinates, since that introduces an unwanted factor of two.

The Clebsch quantities are
\begin{align}
 A={}&-6P_2,\nonumber\\
 B={}&-24P_4+\frac{41}{16}P_2^2,\nonumber\\
 C={}&\frac{576}{13}P_6-\frac{396}{13}P_4P_2
       +\frac{29}{13}P_2^3,\nonumber\\
 D={}&-\frac{62208}{1171}P_8
       +\frac{54864}{1171}P_6P_2
       +\frac{203616}{1171}P_4^2
       -\frac{61287}{1171}P_4P_2^2
       +\frac{13393}{4684}P_2^4,\nonumber\\
 E={}&\frac{41472}{155}P_{10}
       -\frac{4605984}{36301}P_8P_2
       -\frac{106272}{403}P_6P_4
       +\frac{19990440}{471913}P_6P_2^2 \nonumber\\
   &\quad
       +\frac{47719206}{471913}P_4^2P_2
       -\frac{7468023}{471913}P_4P_2^3
       +\frac{10108327}{18876520}P_2^5.
       \label{eq:clebsch}
\end{align}
They are $W(E_6)$-invariant polynomial functions on the six-dimensional source.
Consequently they descend to rational functions, indeed polynomial functions, on
any twist of that source.

If the gamma vector is multiplied by a scalar $\lambda$, then
\[
  (A,B,C,D,E)\mapsto
  (\lambda^2A,\lambda^4B,\lambda^6C,\lambda^8D,\lambda^{10}E).
\]
Thus the moduli point lies in the weighted projective space
\[
  \PP(1,2,3,4,5).
\]

\subsection{From Clebsch invariants to a cubic equation}

A general smooth cubic surface has a Sylvester pentahedral form
\begin{equation}
\begin{cases}
 a_0W_0^3+a_1W_1^3+a_2W_2^3+a_3W_3^3+a_4W_4^3=0,\\
 W_0+W_1+W_2+W_3+W_4=0.
\end{cases}
\label{eq:sylvester}
\end{equation}
Let
\[
  \sigma_i=e_i(a_0,\ldots,a_4),\qquad 1\leq i\leq5.
\]
Given a Clebsch vector $[A:B:C:D:E]$ with $E\neq0$, use
\begin{equation}
 \boxed{
 [\sigma_1,\sigma_2,\sigma_3,\sigma_4,\sigma_5]
 =\left[B,\ D,\ \frac{C^2-AE}{4},\ CE,\ E^2\right].
 }
 \label{eq:sigma-from-clebsch}
\end{equation}
Then form the monic quintic
\begin{equation}
  g(T)=T^5-\sigma_1T^4+\sigma_2T^3-\sigma_3T^2+\sigma_4T-\sigma_5.
  \label{eq:penta-quintic}
\end{equation}
Let
\[
  \cA_5=\QQ[T]/(g),
  \qquad
  \theta=T\bmod g.
\]
Assume $g$ is separable.  Let
\[
  N=\ker\bigl(\Tr_{\cA_5/\QQ}:\cA_5\to\QQ\bigr).
\]
Choose a $\QQ$-basis
\[
  c_0,c_1,c_2,c_3
\]
of $N$ and set
\[
  \ell(X)=c_0X_0+c_1X_1+c_2X_2+c_3X_3
  \in \cA_5[X_0,X_1,X_2,X_3].
\]
The output cubic is
\begin{equation}
 \boxed{
  F(X_0,X_1,X_2,X_3)
  =\Tr_{\cA_5/\QQ}\bigl(\theta\,\ell(X)^3\bigr).
 }
 \label{eq:trace-cubic}
\end{equation}
This is the trace-algebra form of Algorithm A.10.  It avoids computing the
splitting field of the pentahedral quintic.

\section{What the original descent computes}

Let $L/\QQ$ be the splitting field of $f$ and let
$G=\Gal(L/\QQ)$.  The six-dimensional twisted source is
\begin{equation}
  V_\rho=(L\otimes_\QQ \V)^G,
  \label{eq:twisted-source}
\end{equation}
where $G$ acts semilinearly by
\[
  g(a\otimes v)=g(a)\otimes \rho(g)v.
\]
In coordinates, a point of $V_\rho$ is a vector $t\in L^6$ satisfying equations
of the form
\begin{equation}
  M_{\rho(g)}\,g(t)=t,
  \qquad g\in G,
  \label{eq:six-fixed-equation}
\end{equation}
up to the chosen row/column convention.  The original implementation constructs a
basis matrix
\[
  D_\rho\in\operatorname{GL}_6(L)
\]
whose columns form a $\QQ$-basis of this fixed space.

Similarly, the ten-dimensional Coble span has a $G$-twist.  The original
implementation constructs a basis matrix
\[
  B_\rho\in\operatorname{GL}_{10}(L)
\]
for the ten-dimensional semilinear fixed space, descends the thirty Coble cubics
to $\QQ$, searches for a rational point in the descended projective variety, and
then maps the point back to Coble coordinates.

To solve \eqref{eq:six-fixed-equation} and its ten-dimensional analogue by plain
rational linear algebra, one often chooses a $\QQ$-basis
\[
  \omega_1,\ldots,\omega_n
\]
of $L$ and writes every field automorphism $g\in G$ as a matrix in
$\operatorname{GL}_n(\QQ)$.  This is the step we want to avoid.  When
$|G|=[L:\QQ]$ is large, the matrices are large, the field arithmetic is expensive,
and the descent becomes the bottleneck.

\begin{summarybox}
The full splitting field is one model of the torsor, not the only model.  The
algorithm needs the twist
$V_\rho$, not a particular $\QQ$-basis of $L$.  Resolvent algebras provide smaller
models of the same torsor whenever the desired representation is visible inside a
low-degree permutation representation.
\end{summarybox}

\section{First simplification: eliminate the ten-dimensional descent}

The ten-dimensional matrix $B_\rho$ is useful if one wants explicit descended
coordinates on the projective Coble model
\[
  \tM_\rho\subset\PP^9_\QQ.
\]
It is not needed to produce a cubic equation.

Indeed, the trace cubic \eqref{eq:trace-cubic} depends only on the unmarked
Clebsch vector
\[
  [A:B:C:D:E]\in \PP(1,2,3,4,5)(\QQ).
\]
The functions $A,B,C,D,E$ in \eqref{eq:clebsch} are $W(E_6)$-invariant functions
on the Pinkham source $\V$.  Therefore they descend directly to functions
\[
  A_\rho,B_\rho,C_\rho,D_\rho,E_\rho
\]
on the twisted source $V_\rho$.

Thus the production algorithm can follow the shorter route
\[
  u\in V_\rho(\QQ)
  \longmapsto
  \bigl(A_\rho(u),B_\rho(u),C_\rho(u),D_\rho(u),E_\rho(u)\bigr)
  \longmapsto
  \text{trace cubic over }\QQ.
\]
No rational point search in $\PP^9$, no descended Coble cubics, and no
$B_\rho$ are required.

\begin{proposition}[Ten-dimensional descent is unnecessary for equation production]
Assume one can enumerate rational points of the twisted source $V_\rho$ and
evaluate the descended Clebsch functions $A_\rho,B_\rho,C_\rho,D_\rho,E_\rho$ on
those points.  Then one can run the cubic-surface production part of Algorithm 5.1
without constructing the ten-dimensional matrix $B_\rho$ or the descended
projective model $\tM_\rho\subset\PP^9_\QQ$.
\end{proposition}

\begin{proof}
Over $L$, a point of $V_\rho$ becomes an ordinary cusp-parameter vector
$t\in\V$.  The normalized Coble map gives a marked cubic surface, and the power-sum
formulas \eqref{eq:clebsch} give its unmarked Clebsch invariants.  Since these
formulas are $W(E_6)$-invariant, the resulting weighted Clebsch point is fixed by
Galois and is defined over $\QQ$.  Formula \eqref{eq:sigma-from-clebsch} and the
trace construction \eqref{eq:trace-cubic} then produce a cubic equation over
$\QQ$.  The intermediate descended Coble coordinates are only one possible way of
obtaining the same Clebsch point.
\end{proof}

\section{Resolvent algebras attached to $G$-sets}

\subsection{The associated \'etale algebra}

Let $\Omega$ be a finite left $G$-set.  Write $\QQ^\Omega$ for the permutation
module with basis $\{e_\omega:\omega\in\Omega\}$.  Define
\begin{equation}
  \cA_\Omega=(L\otimes_\QQ \QQ^\Omega)^G,
  \label{eq:AOmega}
\end{equation}
where $G$ acts semilinearly on $L$ and by its given permutation action on
$\QQ^\Omega$.

Then $\cA_\Omega$ is an \'etale $\QQ$-algebra of degree $|\Omega|$.  After base
change to $L$, it splits canonically as
\begin{equation}
  L\otimes_\QQ \cA_\Omega\cong L^\Omega.
  \label{eq:AOmega-split}
\end{equation}

If $\Omega$ is transitive and $\Omega\cong G/H$, then
\begin{equation}
  \cA_\Omega\cong L^H.
  \label{eq:transitive-resolvent-field}
\end{equation}
For a general $G$-set decomposed into orbits
\[
  \Omega=\Omega_1\sqcup\cdots\sqcup\Omega_r,
  \qquad
  \Omega_j\cong G/H_j,
\]
one has
\begin{equation}
  \cA_\Omega\cong \prod_{j=1}^r L^{H_j}.
  \label{eq:orbit-product}
\end{equation}

\subsection{Constructing $\cA_\Omega$ from $f$}

Suppose first that $\Omega\cong G/H$ is transitive.  Choose an expression
\[
  \theta_H=\Theta(\alpha_1,\ldots,\alpha_d)\in L
\]
which is fixed by $H$ and whose stabilizer in $G$ is exactly $H$.  Then the
resolvent polynomial
\begin{equation}
  h_H(T)=\prod_{gH\in G/H}\left(T-g(\theta_H)\right)
  \label{eq:resolvent-polynomial}
\end{equation}
has coefficients in $\QQ$ and degree $[G:H]=|\Omega|$.  For a generic enough
choice of $\Theta$, it is separable and
\[
  \QQ[T]/(h_H)\cong L^H\cong\cA_\Omega.
\]

For a nontransitive $G$-set, construct one such resolvent for each orbit and take
the product algebra.  In practice, one chooses $\Theta$ as a suitably generic
linear or low-degree expression in the roots of $f$ and verifies that the computed
resolvent has the expected degree and discriminant nonzero.

\begin{implementationbox}
A list of separate orbit resolvent polynomials gives the underlying product
algebra $\prod_j L^{H_j}$, but some algorithms also need the descended maps induced
by $G$-equivariant linear maps between permutation modules.  Those maps encode the
relative labelling of the different orbits.  Thus the true computational object is
not only the product algebra $\cA_\Omega$, but $\cA_\Omega$ together with the
descended endomorphisms and polynomial functions required below.
\end{implementationbox}

\subsection{Why this avoids the full splitting field}

The full splitting field corresponds to the regular $G$-set $G$ itself:
\[
  (L\otimes_\QQ \QQ^G)^G\cong L.
\]
Using the regular $G$-set therefore gives degree $|G|$ arithmetic, which is exactly
what we are trying to avoid.

A smaller $G$-set $\Omega$ gives degree $|\Omega|$ arithmetic.  The $27$-line
$G$-set always has total degree $27$, often split into smaller orbit factors.  If
$|G|$ is large, this can be dramatically cheaper than computing in $L$.

\section{Constructing the twisted six-dimensional source from a resolvent}

\subsection{General representation-theoretic construction}

Let $V$ be the six-dimensional representation $\V$ of $G$ obtained by restricting
Pinkham's $W(E_6)$ representation along
\[
  \rho:G\to\W.
\]
Choose a finite $G$-set $\Omega$ and a $G$-equivariant surjection
\begin{equation}
  q:\QQ^\Omega\twoheadrightarrow V.
  \label{eq:qOmega}
\end{equation}
Because $\QQ[G]$ is semisimple, $q$ admits a $G$-equivariant section
\begin{equation}
  s:V\hookrightarrow \QQ^\Omega.
  \label{eq:section}
\end{equation}
Set
\begin{equation}
  e=sq\in\End_G(\QQ^\Omega).
  \label{eq:idempotent}
\end{equation}
Then $e$ is a $G$-equivariant idempotent and
\[
  \im(e)\cong V.
\]

Twisting by the $G$-torsor defined by $f$ gives a $\QQ$-linear idempotent
\begin{equation}
  e_\rho:\cA_\Omega\longrightarrow\cA_\Omega.
  \label{eq:descended-idempotent}
\end{equation}
Its image is the desired six-dimensional twisted source.

\begin{proposition}[Resolvent model of the twisted source]
With notation as above, there is a canonical isomorphism
\begin{equation}
  \im(e_\rho)\cong (L\otimes_\QQ V)^G=V_\rho.
  \label{eq:source-as-image}
\end{equation}
In particular, $\im(e_\rho)$ is a six-dimensional $\QQ$-vector space.
\end{proposition}

\begin{proof}
The exact sequence
\[
  0\longrightarrow\ker(q)\longrightarrow \QQ^\Omega
  \xrightarrow{q}V\longrightarrow0
\]
splits $G$-equivariantly by $s$.  Tensoring with $L$ and taking semilinear
$G$-invariants gives
\[
  (L\otimes\QQ^\Omega)^G
  \cong
  (L\otimes\ker(q))^G\oplus (L\otimes V)^G.
\]
Under the identification
\[
  (L\otimes\QQ^\Omega)^G=\cA_\Omega,
\]
the projection onto the second summand is precisely the descended idempotent
$e_\rho$.  Therefore its image is $V_\rho$.
\end{proof}

\subsection{How to compute the idempotent}

Suppose the matrices for the action of $G$ on $\QQ^\Omega$ are $P_g$, and the
matrices for the action on $V$ are $M_g$.  With column vectors, equivariance of
$q$ means
\begin{equation}
  M_g q=qP_g.
  \label{eq:q-equivariance}
\end{equation}
Choose any rational section $s_0$ of $q$, so that $qs_0=I_V$.  Then
\begin{equation}
  s=\frac{1}{|G|}\sum_{g\in G}P_g^{-1}s_0M_g
  \label{eq:averaged-section}
\end{equation}
is a $G$-equivariant section.  Indeed,
\[
  qs=I_V,
  \qquad
  P_hs=sM_h\quad(h\in G).
\]
Then
\[
  e=sq
\]
is the desired rational idempotent on $\QQ^\Omega$.

The descended map $e_\rho$ on $\cA_\Omega$ is the unique $\QQ$-linear map whose
base change to $L$ is the coordinate-wise map induced by $e$ on $L^\Omega$.
Implementation methods include:
\begin{enumerate}
\item constructing the corresponding resolvent correspondences between the orbit
factors $L^{H_i}$ and $L^{H_j}$;
\item computing the map modulo enough good split primes, where
$\cA_\Omega\otimes\mathbf F_p\cong \mathbf F_p^\Omega$, and reconstructing the
rational matrix;
\item in a hybrid test implementation, using the full splitting field only to
calibrate $e_\rho$ once, then replacing it by the lower-degree resolvent model.
\end{enumerate}
The first two methods avoid constructing a $\QQ$-basis of $L$.

\section{Choosing the $G$-set $\Omega$}

\subsection{The root set of $f$}

The cheapest possible choice is the root set of $f$:
\[
  \Omega_f=\{1,\ldots,d\},
\]
with its labelled $G$-action.  Then
\[
  \cA_{\Omega_f}=\QQ[x]/(f)
\]
if $f$ is irreducible; more generally it is the product of the root-algebra factors
corresponding to the $G$-orbits on the roots.

However, this works only if the permutation module $\QQ^{\Omega_f}$ contains the
six-dimensional source representation as a quotient.  A preliminary character test
is
\begin{equation}
  m=\left\langle \chi_V,\chi_{\Omega_f}\right\rangle_G
  =\frac{1}{|G|}\sum_{g\in G}\chi_V(g)\cdot\#\Fix_{\Omega_f}(g).
  \label{eq:character-test}
\end{equation}
If $m=0$, then $\QQ^{\Omega_f}$ cannot contain $V$.

If $m>0$, this is still only a preliminary test, because $V$ may be reducible after
restriction to $G$.  The definitive test is to solve the linear equations
\begin{equation}
  M_g C=CP_g,
  \qquad g\in G,
  \label{eq:intertwiner-equations}
\end{equation}
for matrices
\[
  C\in\Mat_{6\times d}(\QQ),
\]
and check whether some solution has rank $6$.  Such a solution is precisely a
$G$-equivariant surjection
\[
  \QQ^{\Omega_f}\twoheadrightarrow V.
\]

\subsection{The $27$-line set}

The universal safe choice is
\[
  \Omega=\Omega_{27},
\]
the $27$-line set, with $G$-action through $\rho$.  The surjection is the map
$q_{27}$ in \eqref{eq:q27}.  Therefore the rank test is automatic.

The associated algebra is
\[
  \cA_{27}=(L\otimes\QQ^{\Omega_{27}})^G.
\]
It has total degree $27$.  If the image of $G$ has orbits
\[
  \Omega_{27}=\Omega_1\sqcup\cdots\sqcup\Omega_r,
\]
then
\[
  \cA_{27}\cong\prod_{j=1}^r L^{H_j},
  \qquad H_j=\Stab_G(\omega_j).
\]
Even when the product has several factors, the descended idempotent $e_\rho$ is
still defined on the product algebra and has six-dimensional image.

\begin{summarybox}
\textbf{Recommended choice.}
Try the root algebra $\QQ[x]/(f)$ first.  If it contains the six-dimensional source
representation, it is likely the fastest model.  If it does not, use the $27$-line
resolvent algebra.  This gives a uniform degree-$27$ substitute for the full
splitting field.
\end{summarybox}

\subsection{A warning about multiple orbits}

When $\Omega$ has several $G$-orbits, the underlying algebra
\[
  \prod_j L^{H_j}
\]
by itself is not always enough for implementation.  A $G$-equivariant map
\[
  \QQ^\Omega\to V
\]
may mix information from different orbits.  To descend such a map, one needs the
corresponding rational linear maps between the orbit resolvent factors, or an
equivalent method that preserves their common origin from the same $G$-torsor.

Using a single transitive $G$-set whose permutation representation surjects onto
$V$ is technically simpler.  But the $27$-line set may decompose after restriction
to a subgroup $G\subset W(E_6)$, and the product-plus-correspondences model is then
the correct object.

\section{Evaluating invariant functions without $L$}

Let
\[
  F\in \QQ[V]^G
\]
be a $G$-invariant polynomial function on the six-dimensional representation.  In
applications, $F$ will be one of
\[
  \DeltaE,
  \qquad
  A,
  B,
  C,
  D,
  E.
\]
Choose a $G$-equivariant surjection
\[
  q:\QQ^\Omega\twoheadrightarrow V.
\]
Then
\[
  F\circ q\in\QQ[\QQ^\Omega]^G.
\]
Twisting by the torsor defined by $f$ gives an ordinary polynomial function
\begin{equation}
  F_{\rho,\Omega}:\cA_\Omega\longrightarrow\QQ.
  \label{eq:descended-invariant}
\end{equation}
Restricting this function to
\[
  \im(e_\rho)\cong V_\rho
\]
gives the descended function on the twisted source.

Equivalently, if $a\in\im(e_\rho)$ and one base changes to $L$, then $a$ becomes a
vector in $L^\Omega$.  Applying $q$ gives an element of $L\otimes V$, and evaluating
$F$ gives an element of $L$.  Because $F$ is $G$-invariant and $a$ is descended,
that element is fixed by $G$, hence lies in $\QQ$.  This rational number is
$F_{\rho,\Omega}(a)$.

\begin{proposition}[Invariant evaluation]
The values
\[
  \Delta_{E_6,\rho}(a),\quad A_\rho(a),\quad B_\rho(a),\quad C_\rho(a),\quad
  D_\rho(a),\quad E_\rho(a)
\]
for $a\in V_\rho(\QQ)$ can be computed from the resolvent model
$\cA_\Omega$ and the descended idempotent $e_\rho$, without computing a
$\QQ$-basis of the full splitting field $L$.
\end{proposition}

\begin{proof}
The construction above defines each value as the evaluation of a descended
polynomial function on a finite-dimensional $\QQ$-vector space.  The data needed to
write this polynomial are the resolvent algebra $\cA_\Omega$ and the descended
linear maps attached to $q$ and $e$.  The full splitting field appears only after
base change, where it is used to prove that the construction agrees with the
ordinary untwisted invariant $F$.
\end{proof}

\subsection{Practical exact evaluation methods}

There are several implementation strategies.

\begin{enumerate}[label=\textbf{Method \arabic*.}]
\item \textbf{Symbolic descent.}  Fix a $\QQ$-basis of $\cA_\Omega$, write a
general element of $\im(e_\rho)$ in six coordinates, and descend the polynomial
$F\circ q$ to a polynomial in those six coordinates.  This is a one-time
instance-specific computation.

\item \textbf{Modular split-prime evaluation.}  For good primes $p$ at which the
relevant resolvent algebras split and all denominators are invertible, identify
\[
  \cA_\Omega\otimes\mathbf F_p\cong \mathbf F_p^\Omega.
\]
Evaluate $e$, $q$, and the invariant $F$ over $\mathbf F_p$.  Because the result is
invariant under the residual labelling ambiguity, it is a well-defined value modulo
$p$.  Repeat for enough primes and reconstruct the rational value by CRT and
rational reconstruction.  If $\Omega$ has several orbits, the labellings of the
factors must be correlated through the same torsor, or the descended
correspondences must be used.

\item \textbf{Hybrid calibration.}  During development, compute the full splitting
field for small examples, verify the resolvent evaluation against the direct
$L$-evaluation, and then remove the full splitting-field step in production.
\end{enumerate}

The crucial point is that the individual gamma coordinates are not the intended
output of the no-$L$ implementation.  The intended output of the source evaluation
is the invariant vector
\[
  [A:B:C:D:E]\in\PP(1,2,3,4,5)(\QQ).
\]

\section{The improved algorithm}
\label{sec:improved-algorithm}

\subsection{Universal precomputation}

Compute once and store the following data.

\begin{enumerate}[label=\textbf{U\arabic*.}]
\item A concrete rational model of the six-dimensional representation $\V$ of
$\W$, with matrices $M_w$ or a word-evaluation routine for the chosen generators of
$\W$.

\item The $27$-line $W(E_6)$-set and the surjection
\[
  q_{27}:\QQ^{\Omega_{27}}\twoheadrightarrow \V,
  \qquad
  e_\ell\mapsto \ell+K/3.
\]

\item The normalized gamma formulas \eqref{eq:gammahatI} and
\eqref{eq:gammahatII}, the ordering of the forty labels, and the convention sign
mask.

\item The invariant formulas
\[
  \DeltaE=\deltaE^2,
  \qquad
  A,B,C,D,E.
\]
In a direct invariant implementation, one may store these as polynomial functions
on $\V$ rather than as gamma-coordinate routines.

\item The pentahedral reconstruction formulas
\eqref{eq:sigma-from-clebsch}, \eqref{eq:penta-quintic}, and
\eqref{eq:trace-cubic}.
\end{enumerate}

\subsection{Instance-specific computation}

Given $f(x)$, a labelled Galois group $G\leq S_d$, and
$\rho:G\to\W$, proceed as follows.

\begin{enumerate}[label=\textbf{Step \arabic*.}]
\item \textbf{Restrict the source representation.}
Evaluate $\rho(g)$ in the six-dimensional representation $\V$ for generators of
$G$.  This gives matrices $M_g$ for the $G$-action on $V$.

\item \textbf{Choose a low-degree permutation model.}
First test the root set $\Omega_f$ using the intertwiner equations
\eqref{eq:intertwiner-equations}.  If there is a rank-six equivariant map
\[
  \QQ^{\Omega_f}\twoheadrightarrow V,
\]
use $\Omega=\Omega_f$.  Otherwise use $\Omega=\Omega_{27}$ with action through
$\rho$ and $q=q_{27}$.

\item \textbf{Construct the resolvent algebra.}
Construct
\[
  \cA_\Omega=(L\otimes\QQ^\Omega)^G
\]
from $f$ using resolvent polynomials for the orbits of $\Omega$.  Do not construct
a $\QQ$-basis of $L$.

\item \textbf{Construct the descended source.}
Choose a $G$-equivariant surjection
\[
  q:\QQ^\Omega\twoheadrightarrow V
\]
and a $G$-equivariant section $s$.  Form the idempotent
\[
  e=sq.
\]
Descend it to
\[
  e_\rho:\cA_\Omega\to\cA_\Omega.
\]
Compute a $\QQ$-basis
\[
  b_1,\ldots,b_6
\]
for
\[
  \im(e_\rho)\cong V_\rho.
\]
Optionally replace this basis by an LLL-reduced basis for a chosen integral lattice
in $\im(e_\rho)$.

\item \textbf{Enumerate source points.}
Enumerate primitive vectors
\[
  u=(u_1,\ldots,u_6)\in\ZZ^6
\]
by increasing height and set
\[
  a(u)=u_1b_1+\cdots+u_6b_6\in \im(e_\rho).
\]
Primitive enumeration is natural because the Coble and Clebsch constructions are
homogeneous, so scalar multiples usually give the same moduli point.

\item \textbf{Boundary test.}
Evaluate
\[
  \Delta_{E_6,\rho}(a(u)).
\]
If this is zero, reject $u$ and continue.  This is the invariant version of the
condition that no factor in \eqref{eq:deltaE} vanish after base change to $L$.

\item \textbf{Evaluate Clebsch invariants.}
Evaluate
\[
  A=A_\rho(a(u)),\quad
  B=B_\rho(a(u)),\quad
  C=C_\rho(a(u)),\quad
  D=D_\rho(a(u)),\quad
  E=E_\rho(a(u)).
\]
If $[A:B:C:D:E]$ is undefined or $E=0$, reject $u$ and continue.

\item \textbf{Construct the pentahedral quintic.}
Use \eqref{eq:sigma-from-clebsch} to form
\[
  \sigma_1,\ldots,\sigma_5
\]
and then form the quintic $g(T)$ in \eqref{eq:penta-quintic}.  If $g$ is not
separable, reject $u$ and continue.

\item \textbf{Run the trace construction.}
Construct
\[
  \cA_5=\QQ[T]/(g),
\]
choose a basis of the trace-zero subspace, and output the cubic form
\[
  F(X)=\Tr_{\cA_5/\QQ}\bigl(\theta\ell(X)^3\bigr).
\]

\item \textbf{Certify smoothness.}
Check that the resulting cubic surface is smooth.  If it is singular, reject $u$
and continue.  Otherwise return $F$.
\end{enumerate}

\subsection{Language-neutral pseudocode}

\begin{lstlisting}
function CubicSurfaceFromResolvent(f, labelled_G, rho, universal_data):
    # Input:
    #   f: separable polynomial over Q
    #   labelled_G <= S_deg(f), identified with Gal(f)
    #   rho: labelled_G -> W(E6)
    # Output:
    #   smooth cubic surface over Q realizing rho, if a candidate is found

    V_action = restrict_Pinkham_representation(rho)

    if root_permutation_module_has_rank_6_quotient(f, labelled_G, V_action):
        Omega = root_set_of_f
        q = choose_rank_6_intertwiner(Q^Omega -> V)
    else:
        Omega = line_27_G_set_via_rho
        q = q_27

    A_Omega = construct_resolvent_algebra(f, labelled_G, Omega)

    s = equivariant_section(q)
    e = s * q
    e_rho = descend_idempotent_to_resolvent_algebra(e, A_Omega)

    basis = basis_of_image(e_rho)      # six vectors b_1,...,b_6
    basis = optional_LLL_reduce(basis)

    for u in primitive_integer_vectors_by_height(6):
        a = sum(u[i] * basis[i] for i in 1..6)

        if descended_Delta_E(a) == 0:
            continue

        A, B, C, D, E = descended_Clebsch_values(a)
        if weighted_point_is_invalid(A, B, C, D, E):
            continue
        if E == 0:
            continue

        sigmas = [B, D, (C^2 - A*E)/4, C*E, E^2]
        g = T^5 - sigmas[1]*T^4 + sigmas[2]*T^3 \
              - sigmas[3]*T^2 + sigmas[4]*T - sigmas[5]

        if not is_separable(g):
            continue

        F = trace_cubic_from_quintic(g)

        if is_smooth_cubic_surface(F):
            return F

    error("No candidate found in the searched height range")
\end{lstlisting}

\section{Correctness}

\begin{theorem}[Correctness of the improved algorithm]
Assume the universal data use compatible conventions for the $W(E_6)$ action on
$\V$, the $27$ lines, and the normalized Coble gamma coordinates.  Assume also
that the returned candidate lies in the standard open locus where the
Clebsch--pentahedral reconstruction used in Algorithm A.10 recovers the unmarked
cubic surface represented by the Clebsch point.  If the algorithm in
\cref{sec:improved-algorithm} returns a cubic form $F$, then the surface
$X=\{F=0\}\subset\PP^3_\QQ$ is a smooth cubic surface whose Galois action on the
$27$ geometric lines is the prescribed representation $\rho$, up to the initial
marking convention.  The optional final computation of the $27$ lines certifies
this conclusion a posteriori.
\end{theorem}

\begin{proof}
The source space used by the algorithm is
\[
  \im(e_\rho)\cong V_\rho=(L\otimes V)^G
\]
by \cref{eq:source-as-image}.  After base change to $L$, a rational source point
becomes an ordinary vector $t\in V$ satisfying the twisted compatibility condition
with respect to $\rho$.  The nonvanishing of $\Delta_{E_6,\rho}$ means that $t$ avoids
the $E_6$ root hyperplanes, so the six cuspidal points are in general position and
determine a marked smooth cubic surface over $L$.

The descent condition on $t$ says precisely that applying a Galois element
$g\in G$ to the marked surface is the same as changing the marking by
$\rho(g)\in W(E_6)$.  Hence the associated Galois action on the marked line
configuration is the one prescribed by $\rho$.

The Clebsch quantities $A,B,C,D,E$ are $W(E_6)$-invariant.  Therefore their values
on the twisted source point are fixed by $G$ and lie in $\QQ$.  These values are
the unmarked moduli invariants of the cubic surface obtained over $L$.

The pentahedral and trace construction from
\eqref{eq:sigma-from-clebsch}--\eqref{eq:trace-cubic} constructs a cubic surface
over $\QQ$ with the same Clebsch invariants.  On the open locus where the
pentahedral construction is valid and the final smoothness check passes, this
surface is the desired unmarked cubic surface.  Since the marked object over $L$
descends with cocycle $\rho$, its line action is the prescribed one, up to the
fixed convention identifying markings with the standard $E_6$ configuration.
\end{proof}

\begin{proposition}[Termination of exhaustive search]
Assume the enumeration of primitive integer vectors in the chosen source lattice is
exhaustive by height.  Then, outside the proper failure loci of the Coble/Clebsch
and pentahedral constructions, the search eventually finds a valid point.
\end{proposition}

\begin{proof}
Over $L$, the source map
\[
  \Aff^6\dashrightarrow \tcM
\]
coming from six points on the cuspidal cubic is dominant.  Twisting by the cocycle
$\rho$ gives a dominant rational map
\[
  V_\rho\dashrightarrow \tcM_\rho.
\]
The source $V_\rho$ is a six-dimensional vector space over $\QQ$, so its rational
points are Zariski dense.  The boundary $\DeltaE=0$, the locus $E=0$, the
inseparability locus of the pentahedral quintic, and the singular-output locus are
proper closed conditions on the source, provided one is not in a degenerate
component of the construction.  Therefore rational points avoiding them are
Zariski dense.  Exhaustive primitive enumeration must eventually meet this open
set.  The argument gives no useful universal height bound.
\end{proof}

\section{Certification and debugging checks}

A robust implementation should perform the following checks.

\subsection{Universal checks}

\begin{enumerate}
\item Verify that the stored matrices for Pinkham's six-dimensional representation
satisfy the Coxeter relations for $W(E_6)$.

\item Verify that the $27$-line permutation action agrees with the action on the
Picard lattice and that $q_{27}$ is $W(E_6)$-equivariant:
\[
  M_wq_{27}=q_{27}P_w.
\]

\item Verify the normalized gamma equivariance
\[
  \Gamhat(M_wt)=\widehat P_w\Gamhat(t)
\]
for the chosen generators and several random rational regular vectors $t$.

\item Verify that the power-sum formulas \eqref{eq:clebsch} are invariant under the
stored $W(E_6)$ action.

\item Verify that the implementation sums over forty gamma labels, not over eighty
signed labels.
\end{enumerate}

\subsection{Instance checks}

\begin{enumerate}
\item Verify that the labelled Galois group of $f$ has the declared order and acts
on the roots in the declared way.

\item If the root algebra is used, verify that the chosen intertwiner
\[
  \QQ^{\Omega_f}\to V
\]
has rank $6$ and satisfies \eqref{eq:intertwiner-equations} for the generators of
$G$.

\item If the $27$-line algebra is used, verify that the orbit resolvents have the
expected degrees and are separable.

\item Verify that the descended idempotent satisfies
\[
  e_\rho^2=e_\rho
\]
and that
\[
  \dim_\QQ\im(e_\rho)=6.
\]

\item For candidates, verify that
\[
  \Delta_{E_6,\rho}(a)\neq0
\]
before evaluating the higher-degree Clebsch functions.

\item Verify that the computed Clebsch vector is not the zero weighted point and
that $E\neq0$ before constructing the pentahedral quintic.

\item Verify separability of the pentahedral quintic and smoothness of the final
cubic surface.

\item As an independent final certification, compute the $27$ lines of the output
surface and compare the resulting Galois action with the prescribed representation
$\rho$.
\end{enumerate}

\section{A staged implementation plan}

The full resolvent implementation is the clean mathematical endpoint, but it need
not be implemented all at once.

\subsection{Stage 1: remove $B_\rho$}

Keep the existing full-splitting-field construction of the six-dimensional matrix
$D_\rho$, but delete the ten-dimensional point-search machinery.  For each
primitive vector $u\in\ZZ^6$, form
\[
  t=D_\rho u\in L^6,
\]
check $\deltaE(t)\neq0$, evaluate the normalized gamma formulas, compute
$A,B,C,D,E$, coerce these invariant values to $\QQ$, and run the pentahedral trace
construction.  This tests the idea that the ten-dimensional descent is unnecessary.

\subsection{Stage 2: try the root algebra}

Before building any $27$-line resolvents, test whether the root permutation module
of $f$ contains $V$ as a quotient.  If so, the algebra
\[
  \QQ[x]/(f)
\]
may already support the source descent.  This is the cheapest possible no-full-$L$
model.

\subsection{Stage 3: use the $27$-line resolvent}

If the root algebra does not contain $V$, construct the $27$-line resolvent algebra
and the descended idempotent $e_\rho$.  This gives a uniform degree-$27$ model of
the twisted source.  The main technical work is implementing the descended linear
maps and invariant evaluators without losing the relative labelling of orbit
factors.

\begin{summarybox}
\textbf{Practical expectation.}
Stage 1 should already be faster than the original ten-dimensional search.  Stage 2
can be much faster when the given polynomial's root action already contains the
six-dimensional representation.  Stage 3 is the robust general method for avoiding
a full explicit $\QQ$-basis of the splitting field.
\end{summarybox}

\section{What remains expensive}

The resolvent approach removes the need to do linear algebra in a vector space of
dimension $[L:\QQ]=|G|$, but it does not make the arithmetic trivial.  The following
steps may still be nontrivial.

\begin{itemize}
\item Constructing resolvent polynomials can be expensive if the chosen resolvent
has large degree, large coefficients, or difficult stabilizer.

\item If $\Omega$ has several orbits, the descended idempotent and invariant
functions require correspondences between orbit factors, not merely independent
polynomial fields.

\item The Clebsch functions have high degree in the source coordinates.  Their
values may have large numerators and denominators.  Modular evaluation and rational
reconstruction may be preferable to direct symbolic expansion.

\item The algorithm gives a dense supply of rational points but no small height
bound.  LLL reduction and good source lattices remain important.
\end{itemize}

These are implementation costs inside degree $d$ or degree $27$ \'etale algebras,
not costs inside the full splitting field of degree $|G|$.

\section{Minimal final form of the algorithm}

The mathematically compressed version is the following.

\begin{enumerate}[label=\textbf{A\arabic*.}]
\item From $f$ and the labelled group $G\leq S_{\deg f}$, regard
$\Spec L\to\Spec\QQ$ as a $G$-torsor, without constructing $L$ as a
$\QQ$-vector space.

\item Restrict Pinkham's six-dimensional $W(E_6)$ representation along
$\rho:G\to W(E_6)$.

\item Choose a finite $G$-set $\Omega$ and a $G$-equivariant surjection
$\QQ^\Omega\twoheadrightarrow V$.  Use the root set of $f$ if it passes the rank
six intertwiner test; otherwise use the $27$-line set.

\item Construct the associated resolvent algebra
\[
  \cA_\Omega=(L\otimes\QQ^\Omega)^G
\]
and the descended idempotent $e_\rho$ whose image is the twisted source
$V_\rho$.

\item Enumerate rational points of
\[
  V_\rho=\im(e_\rho).
\]

\item For each candidate, evaluate the descended invariant functions
\[
  \Delta_{E_6,\rho},\quad
  A_\rho,B_\rho,C_\rho,D_\rho,E_\rho.
\]
Reject the candidate if it lies on the boundary or outside the valid pentahedral
locus.

\item Convert the Clebsch vector to a pentahedral quintic and then to a cubic form
by the trace construction.

\item Check smoothness and, for certification, verify the induced Galois action on
the $27$ lines.
\end{enumerate}

\begin{thebibliography}{9}

\bibitem{EJ}
A.-S. Elsenhans and J. Jahnel,
\emph{Moduli spaces and the inverse Galois problem for cubic surfaces},
Trans. Amer. Math. Soc. \textbf{367} (2015), no.~11, 7837--7861;
\href{https://arxiv.org/abs/1209.5591}{arXiv:1209.5591}.

\bibitem{DR-PR}
A. Duncan and Z. Reichstein,
\emph{Pseudo-reflection groups and essential dimension},
arXiv:1307.5724; see especially Lemma~6.1.

\bibitem{DR-versality}
A. Duncan and Z. Reichstein,
\emph{Versality of algebraic group actions and rational points on twisted varieties},
J. Algebraic Geom. \textbf{24} (2015), 499--530;
\href{https://arxiv.org/abs/1109.6093}{arXiv:1109.6093}.

\bibitem{Pinkham}
H. Pinkham,
\emph{R\'esolution simultan\'ee de points doubles rationnels},
in \emph{S\'eminaire sur les singularit\'es des surfaces},
Lecture Notes in Math. \textbf{777}, Springer, 1980, 179--203.

\bibitem{Cohen}
H. Cohen,
\emph{A Course in Computational Algebraic Number Theory},
Graduate Texts in Mathematics \textbf{138}, Springer, 1993.

\bibitem{Stauduhar}
R. P. Stauduhar,
\emph{The determination of Galois groups},
Math. Comp. \textbf{27} (1973), 981--996.

\bibitem{EJcode}
J. Jahnel,
\emph{c9\_example.m: Magma example for Algorithm 5.1},
authors' supplementary webpage,
\url{https://wwwuser.gwdguser.de/~jjahnel/Arbeiten/Kub_Fl/c9_example.m}.

\end{thebibliography}

\end{document}
